\section{圆锥曲线相似的条件}
\markboth{89. 圆锥曲线相似的条件}{89. 圆锥曲线相似的条件}

设曲线 $C$ 是一条圆锥曲线，那么它的极坐标方程总可以写成形式
\[
  \rho = \frac{p}{1 - \varepsilon\cos\theta},
\]
式中 $\varepsilon$ 是曲线 $C$ 的离心率，$p$ 是一个不等于零的常数。取极点 $O$ 为位似中心，作位似变换，把曲线 $C$ 变到 $C_1$，由第 88 节知道 $C_1$ 的方程是
\[
  \frac{\rho}{k} = \frac{p}{1 - \varepsilon\cos\theta},
\]
即
\[
  \rho = \frac{kp}{1 - \varepsilon\cos\theta}.
\]
记 $kp = p_1$，化为
\[
  \rho = \frac{p_1}{1 - \varepsilon\cos\theta},
\]
这说明 $C_1$ 仍是离心率为 $\varepsilon$ 的圆锥曲线。

反过来，设 $C$ 和 $C_1$ 是任意两条具有相同离心率 $\varepsilon$ 的圆锥曲线，那么可把它们移动到适当位置，使 $C$ 的极坐标方程是 $\rho = \frac{p}{1 - \varepsilon\cos\theta}$，同时 $C_1$ 的方程是 $\rho = \frac{p_1}{1 - \varepsilon\cos\theta}$。取极点为位似中心，$k = \frac{p_1}{p}$，作位似变换，就能把曲线 $C$ 变到 $C_1$，因而这两条曲线相似。

综上所述，可知要使两条圆锥曲线相似，必须而且只须它们的离心率相同。

作为特例，由于所有的抛物线的离心率都等于 $1$，所以立刻推出所有抛物线都彼此相似。
